% !TeX program = xelatex
% ============================================================
%  RAPPORT DE PROJET – AUTOMATISME INDUSTRIEL
%  Commande d'un Robot KUKA par API en Langage Ladder
%  Mahmoudi Assaad – Génie Électrique / Automatique Industrielle
%  Année Universitaire 2025–2026
% ============================================================

\documentclass[12pt, a4paper]{report}

% ── Encodage, langue & police ────────────────────────────────
\usepackage{iftex}
\ifPDFTeX
  \usepackage[utf8]{inputenc}
  \usepackage[T1]{fontenc}
  \usepackage{lmodern}
\else
  \usepackage{fontspec}
  \setmainfont{Latin Modern Roman}
  \setsansfont{Latin Modern Sans}
  \setmonofont{Latin Modern Mono}
\fi
\usepackage[french]{babel}

% ── Mise en page ─────────────────────────────────────────────
\usepackage[
  a4paper,
  top=2.5cm, bottom=2.5cm,
  left=3cm, right=2.5cm,
  headheight=16pt
]{geometry}
\usepackage{setspace}
\onehalfspacing

% ── Typographie ──────────────────────────────────────────────
\usepackage{microtype}
\usepackage{csquotes}
\raggedbottom
\emergencystretch=3em
\hfuzz=10000pt
\vfuzz=10000pt
\hbadness=10000
\vbadness=10000

% ── Couleurs ─────────────────────────────────────────────────
\usepackage[dvipsnames, table]{xcolor}
\definecolor{bleuFonce}{RGB}{31,  78, 121}
\definecolor{bleuMoyen}{RGB}{46, 117, 182}
\definecolor{bleuClair}{RGB}{214, 228, 240}
\definecolor{orange}{RGB}{197,  90,  17}
\definecolor{grisClaire}{RGB}{242, 242, 242}
\definecolor{vertCode}{RGB}{0,  128,   0}
\definecolor{rougeCode}{RGB}{163,  21,  21}
\definecolor{codeBg}{RGB}{248, 248, 248}
\definecolor{vertEtape}{RGB}{0, 153, 76}
\definecolor{rougeAlerte}{RGB}{192, 0, 0}

% ── Graphiques & TikZ ────────────────────────────────────────
\usepackage{graphicx}
\usepackage{float}
\usepackage{caption}
\usepackage{subcaption}
\usepackage{tikz}
\usetikzlibrary{shapes.geometric, arrows.meta, positioning,
                fit, backgrounds, shadows, calc, decorations.pathreplacing}

% ── Tableaux ─────────────────────────────────────────────────
\usepackage{booktabs}
\usepackage{tabularx}
\usepackage{longtable}
\usepackage{multirow}
\usepackage{array}
\usepackage{colortbl}
\newcolumntype{C}[1]{>{\centering\arraybackslash}p{#1}}
\newcolumntype{L}[1]{>{\raggedright\arraybackslash}p{#1}}
\setlength{\tabcolsep}{4pt}

% ── Code source ──────────────────────────────────────────────
\usepackage{listings}
\lstdefinestyle{ladder}{
  backgroundcolor=\color{codeBg},
  basicstyle=\ttfamily\footnotesize,
  breaklines=true,
  frame=single,
  rulecolor=\color{bleuMoyen},
  framesep=6pt,
  xleftmargin=10pt,
  xrightmargin=10pt,
  numbers=left,
  numberstyle=\tiny\color{gray},
  numbersep=8pt,
  tabsize=2,
  showstringspaces=false,
  commentstyle=\color{vertCode}\itshape,
  keywordstyle=\color{bleuFonce}\bfseries,
}

% ── En-têtes & pieds de page ─────────────────────────────────
\usepackage{fancyhdr}
\pagestyle{fancy}
\fancyhf{}
\fancyhead[L]{\small\color{bleuMoyen}\leftmark}
\fancyhead[R]{\small\color{bleuMoyen}\rightmark}
\fancyfoot[C]{\small\color{bleuMoyen}\thepage}
\fancyfoot[L]{\small\color{gray}Mahmoudi Assaad --- 2025/2026}
\renewcommand{\headrulewidth}{0.4pt}
\renewcommand{\footrulewidth}{0.4pt}

% ── Boîtes ───────────────────────────────────────────────────
\usepackage{tcolorbox}
\tcbuselibrary{skins, breakable}

\newtcolorbox{infobox}[1][]{
  colback=bleuClair, colframe=bleuMoyen,
  fonttitle=\bfseries\color{bleuFonce},
  title={#1}, arc=4pt, boxrule=1pt,
  left=8pt, right=8pt, top=6pt, bottom=6pt,
}
\newtcolorbox{alertbox}[1][]{
  colback=orange!10, colframe=orange,
  fonttitle=\bfseries\color{orange},
  title={#1}, arc=4pt, boxrule=1pt,
}

% ── Liens hypertexte ─────────────────────────────────────────
\usepackage[
  hypertexnames=false,
  colorlinks=true,
  linkcolor=bleuFonce,
  citecolor=bleuMoyen,
  urlcolor=orange,
  pdftitle={Rapport - Commande Robot KUKA par Langage Ladder},
  pdfauthor={Mahmoudi Assaad}
]{hyperref}

% ── Mathématiques ────────────────────────────────────────────
\usepackage{amsmath, amssymb}
\usepackage{siunitx}
\sisetup{locale=FR, per-mode=symbol}

% ── Listes & bibliographie ───────────────────────────────────
\usepackage{enumitem}
\usepackage[
  backend=biber,
  style=ieee,
  sorting=none
]{biblatex}
\addbibresource{references.bib}

% ── Macros ───────────────────────────────────────────────────
\newcommand{\plc}{PLC}
\newcommand{\api}{API}
\newcommand{\kuka}{KUKA}
\newcommand{\krl}{KRL}
\newcommand{\profinet}{PROFINET}
\newcommand{\tia}{TIA Portal}
\newcommand{\siemens}{Siemens S7-1200}

% ============================================================
\begin{document}
% ============================================================

% ── Page de garde ────────────────────────────────────────────
\begin{titlepage}
  \centering
  \vspace*{1.5cm}

  \begin{tikzpicture}
    \fill[bleuFonce] (0,0) rectangle (\textwidth, 1.4cm);
    \node[white, font=\large\bfseries, text width=0.95\textwidth, align=center]
      at (0.5\textwidth, 0.7cm)
      {D\'EPARTEMENT DE G\'ENIE \'ELECTRIQUE --- AUTOMATIQUE INDUSTRIELLE};
  \end{tikzpicture}

  \vspace{1.5cm}
  {\color{gray}\large Rapport de Projet de Fin d'\'Etudes}\\[0.5cm]
  \rule{0.6\textwidth}{0.6pt}
  \vspace{1cm}

  {\huge\bfseries\color{bleuFonce}
    Conception et R\'ealisation d'un\\[0.3cm]
    Syst\`eme Automatis\'e de Tri\\[0.3cm]
    par Robot Industriel \kuka{}}

  \vspace{0.8cm}
  {\Large\color{bleuMoyen}
    Commande par \api{} Siemens S7-1200\\
    Programmation en Langage Ladder}

  \vspace{1.2cm}
  \rule{0.6\textwidth}{0.6pt}
  \vspace{1.2cm}

  \begin{tcolorbox}[
    colback=bleuClair, colframe=bleuMoyen,
    width=0.75\textwidth, arc=6pt,
    left=12pt, right=12pt, top=10pt, bottom=10pt
  ]
    \centering
    {\large\bfseries R\'ealis\'e par :}\\[0.4cm]
    {\Large\color{bleuFonce}\textbf{Mahmoudi Assaad}}\\[0.5cm]
    {\normalsize Niveau : G\'enie \'Electrique --- Automatique Industrielle}\\[0.3cm]
    {\normalsize Ann\'ee Universitaire : \textbf{2025 -- 2026}}
  \end{tcolorbox}

  \vfill
  \begin{tikzpicture}
    \fill[bleuFonce] (0,0) rectangle (\textwidth, 0.9cm);
    \node[white, font=\normalsize] at (0.5\textwidth, 0.45cm)
      {Automatisme Industriel --- Robotique --- Programmation Ladder};
  \end{tikzpicture}
\end{titlepage}

\cleardoublepage

% ── D\'edicace ─────────────────────────────────────────────────
\chapter*{D\'edicace}
\thispagestyle{empty}
\vspace*{4cm}
\begin{flushright}
  \itshape\large
  \`A ma famille, pour son soutien ind\'efectible\\
  tout au long de ce parcours universitaire.\\[1cm]
  \`A mes enseignants et encadrants,\\
  pour la qualit\'e de leur transmission du savoir.\\[1cm]
  \`A tous ceux qui ont cru en mes capacit\'es.
\end{flushright}
\cleardoublepage

% ── Remerciements ────────────────────────────────────────────
\chapter*{Remerciements}
\addcontentsline{toc}{chapter}{Remerciements}

Je tiens \`a exprimer ma profonde gratitude \`a l'ensemble des personnes qui ont contribu\'e,
de pr\`es ou de loin, \`a la r\'ealisation de ce projet d'automatisme industriel.

Mes sinc\`eres remerciements vont en premier lieu \`a mes encadrants p\'edagogiques du
d\'epartement G\'enie \'Electrique, dont les conseils avis\'es, la disponibilit\'e et la rigueur
scientifique ont \'et\'e d\'eterminants pour la r\'eussite de ce travail.

Je remercie \'egalement les responsables du laboratoire d'automatisme industriel pour
l'acc\`es aux \'equipements de simulation et aux ressources documentaires n\'ecessaires.

\cleardoublepage

% ── R\'esum\'e ───────────────────────────────────────────────────
\chapter*{R\'esum\'e}
\addcontentsline{toc}{chapter}{R\'esum\'e}

\begin{infobox}[R\'esum\'e --- Fran\c{c}ais]
Ce rapport pr\'esente la conception et la r\'ealisation compl\`ete d'un syst\`eme automatis\'e de
tri et de manipulation de pi\`eces industrielles. Le syst\`eme int\`egre un robot industriel
\kuka{} \`a 6 axes command\'e par un automate programmable \siemens{}, programm\'e
en langage \textbf{Ladder} sous l'environnement \tia{}.

La communication entre l'automate et le contr\^oleur robot KRC4 est assur\'ee par le
protocole \profinet{} IO. Le document couvre l'analyse fonctionnelle (GRAFCET),
la programmation Ladder compl\`ete, les sch\'emas \'electriques et la gestion de la
s\'ecurit\'e industrielle (PLd/SIL2).

\textbf{Mots-cl\'es :} Robot \kuka{}, \api{} \siemens{}, Langage Ladder,
\profinet{}, GRAFCET, Automatisme industriel, S\'ecurit\'e fonctionnelle.
\end{infobox}

\vspace{1cm}

\begin{infobox}[Abstract --- English]
This report presents the complete design and implementation of an automated industrial
part sorting and handling system. The system integrates a 6-axis \kuka{} industrial
robot controlled by a Siemens S7-1200 PLC, programmed in \textbf{Ladder language}
under TIA Portal.

Communication between the PLC and the KRC4 robot controller is handled via
\profinet{} IO. The document covers functional analysis (GRAFCET), complete Ladder
programming, electrical schematics and industrial safety management (PLd/SIL2).

\textbf{Keywords:} \kuka{} Robot, Siemens S7-1200 PLC, Ladder Language,
\profinet{}, GRAFCET, Industrial Automation, Functional Safety.
\end{infobox}

\cleardoublepage

% ── Tables ───────────────────────────────────────────────────
\tableofcontents
\cleardoublepage
\listoffigures
\addcontentsline{toc}{chapter}{Liste des figures}
\cleardoublepage
\listoftables
\addcontentsline{toc}{chapter}{Liste des tableaux}
\cleardoublepage

% ── Abr\'eviations ─────────────────────────────────────────────
\chapter*{Liste des Abr\'eviations}
\addcontentsline{toc}{chapter}{Liste des abr\'eviations}

\begin{tabular}{@{}lL{10cm}@{}}
  \toprule
  \textbf{Abr\'eviation} & \textbf{Signification} \\
  \midrule
  \api{}    & Automate Programmable Industriel \\
  \plc{}    & Programmable Logic Controller \\
  \kuka{}   & Keller und Knappich Augsburg \\
  \krl{}    & KUKA Robot Language \\
  KRC4      & KUKA Robot Controller 4\textsuperscript{e} g\'en\'eration \\
  \profinet & Process Field Network (Ethernet industriel) \\
  \tia      & Totally Integrated Automation Portal (Siemens) \\
  GRAFCET   & Graphe Fonctionnel de Commande \'Etape/Transition \\
  IEC       & International Electrotechnical Commission \\
  SIL       & Safety Integrity Level \\
  PL        & Performance Level \\
  TOR       & Tout Ou Rien \\
  NO / NF   & Normalement Ouvert / Normalement Ferm\'e \\
  FDC       & Fin De Course \\
  EV        & \'Electrovalve \\
  \bottomrule
\end{tabular}

\cleardoublepage

% ============================================================
%  CHAPITRE 1
% ============================================================
\chapter{Introduction G\'en\'erale}

\section{Contexte industriel}

L'automatisation industrielle constitue aujourd'hui l'un des piliers fondamentaux
des cha\^ines de production modernes. Face aux exigences croissantes de productivit\'e,
de qualit\'e et de s\'ecurit\'e, les entreprises industrielles font appel \`a des syst\`emes
automatis\'es int\'egrant des robots, des automates programmables et des r\'eseaux de
communication industriels.

La quatri\`eme r\'evolution industrielle --- communément d\'esign\'ee sous le terme
\emph{Industrie 4.0} --- acc\'el\`ere cette transformation en int\'egrant des technologies
num\'eriques telles que l'Internet des Objets industriel (\emph{IIoT}),
l'intelligence artificielle et les syst\`emes cyber-physiques.

\section{Probl\'ematique}

\begin{alertbox}[Probl\'ematique du projet]
  Comment concevoir et programmer un syst\`eme automatis\'e de tri de pi\`eces int\'egrant
  un robot industriel \kuka{} command\'e par un \api{} \siemens{} en langage Ladder,
  en assurant la s\'ecurit\'e des op\'erateurs et la fiabilit\'e du cycle de production ?
\end{alertbox}

\section{Objectifs du projet}

\begin{enumerate}[label=\textbf{\arabic*.}, leftmargin=*, itemsep=4pt]
  \item Concevoir l'architecture mat\'erielle et logicielle du syst\`eme automatis\'e
  \item Programmer l'\api{} en langage Ladder sous \tia{}
  \item Assurer la communication PLC -- robot \kuka{} via \profinet{}
  \item Mod\'eliser le comportement s\'equentiel par GRAFCET
  \item Impl\'ementer les fonctions de s\'ecurit\'e industrielle (PLd/SIL2)
  \item Simuler et valider le fonctionnement global du syst\`eme
\end{enumerate}

% ============================================================
%  CHAPITRE 2 — ROBOT KUKA
% ============================================================
\chapter{Pr\'esentation du Robot KUKA}
\label{chap:kuka}

\section{Architecture m\'ecanique du robot 6 axes}

\begin{table}[H]
  \centering
  \caption{Description des axes du robot \kuka{} 6 axes}
  \label{tab:axes}
  \rowcolors{2}{grisClaire}{white}
  \begin{tabular}{C{1.2cm} C{3cm} L{9cm}}
    \toprule
    \rowcolor{bleuMoyen}
    \color{white}\textbf{Axe} &
    \color{white}\textbf{D\'esignation} &
    \color{white}\textbf{Description du mouvement} \\
    \midrule
    A1 & Base         & Rotation de la base autour de l'axe vertical \\
    A2 & Bras inf.    & Balancement du bras inf\'erieur \\
    A3 & Bras sup.    & \'El\'evation du bras sup\'erieur \\
    A4 & Poignet rot. & Rotation du poignet autour de l'axe du bras \\
    A5 & Poignet inc. & Inclinaison du poignet \\
    A6 & Bride        & Rotation finale de la bride porte-outil \\
    \bottomrule
  \end{tabular}
\end{table}

\section{Caract\'eristiques techniques}

\begin{table}[H]
  \centering
  \caption{Caract\'eristiques techniques du robot \kuka{} KR~6 R900}
  \label{tab:carac_kuka}
  \rowcolors{2}{grisClaire}{white}
  \begin{tabular}{L{5.5cm} L{8.5cm}}
    \toprule
    \rowcolor{bleuMoyen}
    \color{white}\textbf{Param\`etre} &
    \color{white}\textbf{Valeur / Sp\'ecification} \\
    \midrule
    Nombre d'axes            & 6 (robot articul\'e s\'erie KR) \\
    Charge utile nominale    & \SI{6}{\kilogram} \\
    Rayon d'action           & \SI{900}{\milli\meter} \\
    R\'ep\'etabilit\'e       & $\pm$\,\SI{0.05}{\milli\meter} \\
    Masse robot              & \SI{52}{\kilogram} \\
    Type de contr\^oleur     & KRC4 compact \\
    Langage de programmation & \krl{} (KUKA Robot Language) \\
    Communication            & \profinet{}, EtherNet/IP \\
    Alimentation             & \SI{400}{\volt} AC / \SI{50}{\hertz} \\
    Degr\'e de protection    & IP54 \\
    Temp\'erature de service & \SI{10}{\degreeCelsius} \`a \SI{55}{\degreeCelsius} \\
    \bottomrule
  \end{tabular}
\end{table}

\section{Sch\'ema du robot 6 axes}

\begin{figure}[H]
  \centering
  \begin{tikzpicture}[scale=1.0]
    \fill[bleuFonce!80] (-1.2,0) rectangle (1.2,0.5);
    \node[white, font=\small\bfseries] at (0,0.25) {SOCLE};

    \fill[bleuMoyen] (-0.7,0.5) rectangle (0.7,1.5);
    \draw[white, very thick] (0,0.5) -- (0,1.5);
    \node[white, font=\footnotesize\bfseries] at (-1.5,1.0) {A1};
    \draw[-{Stealth}, orange, thick] (0.9,1.0) arc (0:270:0.35);
    \node[orange, font=\footnotesize\bfseries] at (1.5,0.75) {Rot.};

    \fill[bleuMoyen!70] (-0.4,1.5) rectangle (0.4,3.2);
    \fill[bleuFonce] (-0.5,1.4) rectangle (0.5,1.65);
    \node[white, font=\footnotesize\bfseries] at (-1.5,2.3) {A2};
    \draw[-{Stealth}, orange, thick] (0.55,1.55) arc (-90:90:0.4);
    \node[orange, font=\footnotesize\bfseries] at (1.4,1.85) {Bascule};

    \fill[bleuMoyen!50] (-0.35,3.2) rectangle (0.35,4.6);
    \fill[bleuFonce] (-0.45,3.1) rectangle (0.45,3.35);
    \node[white, font=\footnotesize\bfseries] at (-1.5,3.9) {A3};
    \draw[-{Stealth}, orange, thick] (0.5,3.22) arc (-90:90:0.35);
    \node[orange, font=\footnotesize\bfseries] at (1.4,3.5) {Elev.};

    \fill[bleuFonce!60] (-0.3,4.6) rectangle (0.3,5.2);
    \fill[bleuMoyen!40] (-0.25,5.2) rectangle (0.25,5.7);
    \node[font=\footnotesize\bfseries, bleuFonce] at (-1.5,4.9) {A4};
    \node[font=\footnotesize\bfseries, bleuFonce] at (-1.5,5.45) {A5};
    \draw[-{Stealth}, orange, thick] (0.38,4.9) arc (0:180:0.18);
    \draw[-{Stealth}, orange, thick] (0.35,5.45) arc (-90:90:0.22);

    \fill[orange!80] (-0.22,5.7) rectangle (0.22,6.0);
    \fill[orange] (-0.15,6.0) rectangle (0.15,6.25);
    \node[font=\footnotesize\bfseries, orange] at (-1.5,5.9) {A6};
    \draw[-{Stealth}, bleuFonce, thick] (0.3,5.85) arc (0:270:0.22);

    \draw[gray, dashed, thin] (-0.7,0.25) -- (-1.2,0.25);
    \draw[gray, dashed, thin] (-0.7,2.3)  -- (-1.2,2.3);
    \draw[gray, dashed, thin] (-0.45,3.9) -- (-1.2,3.9);
    \draw[gray, dashed, thin] (-0.38,4.9) -- (-1.2,4.9);
    \draw[gray, dashed, thin] (-0.35,5.45)-- (-1.2,5.45);
    \draw[gray, dashed, thin] (-0.25,5.9) -- (-1.2,5.9);

    \node[font=\footnotesize, align=left] at (5.5,3.5)
      {\textbf{L\'egende :}\\[4pt]
       \textcolor{bleuFonce}{$\blacksquare$} Structure m\'ecanique\\
       \textcolor{orange}{$\blacksquare$} Articulations / outil\\
       \textcolor{orange}{$\rightarrow$} Sens de rotation};
  \end{tikzpicture}
  \caption{Sch\'ema simplifi\'e du robot \kuka{} 6 axes}
  \label{fig:robot6axes}
\end{figure}

% ============================================================
%  CHAPITRE 3 — AUTOMATE PROGRAMMABLE
% ============================================================
\chapter{Pr\'esentation de l'Automate Programmable}
\label{chap:plc}

\section{Siemens S7-1200 : sp\'ecifications}

\begin{table}[H]
  \centering
  \caption{Sp\'ecifications techniques du \siemens{}}
  \label{tab:s71200}
  \rowcolors{2}{grisClaire}{white}
  \begin{tabular}{L{5.5cm} L{8.5cm}}
    \toprule
    \rowcolor{bleuMoyen}
    \color{white}\textbf{Param\`etre} &
    \color{white}\textbf{Sp\'ecification} \\
    \midrule
    M\'emoire programme        & 50 \`a \SI{150}{\kilo\byte} \\
    E/S num\'eriques int\'egr\'ees & 14 DI / 10 DO (CPU 1214C) \\
    Entr\'ees analogiques      & 2 AI int\'egr\'ees \\
    Port Ethernet              & 1 port RJ45 \profinet{} int\'egr\'e \\
    Protocoles                 & \profinet{} IO, Modbus TCP, TCP/IP \\
    Temps de cycle             & \SI{0.085}{\micro\second}/instruction \\
    Alimentation               & \SI{24}{\volt} DC \\
    Modules d'extension        & Jusqu'\`a 8 modules SM/SB \\
    \bottomrule
  \end{tabular}
\end{table}

\section{Sch\'ema du cycle de scrutation}

\begin{figure}[H]
  \centering
  \begin{tikzpicture}[
    bloc/.style={
      draw=bleuMoyen, fill=bleuClair, rounded corners=6pt,
      minimum width=4.5cm, minimum height=1.2cm,
      font=\bfseries\normalsize, align=center,
      drop shadow={shadow xshift=2pt, shadow yshift=-2pt, opacity=0.3}
    },
    arr/.style={-{Stealth[length=8pt]}, very thick, bleuFonce}
  ]
    \node[bloc, fill=bleuFonce, text=white] (init) at (0,0)
      {Initialisation\\CPU};
    \node[bloc] (lect) at (0,-2.2)
      {\textcolor{bleuFonce}{1.} Lecture\\des entr\'ees};
    \node[bloc] (exec) at (0,-4.4)
      {\textcolor{bleuFonce}{2.} Ex\'ecution\\du programme};
    \node[bloc] (ecri) at (0,-6.6)
      {\textcolor{bleuFonce}{3.} \'Ecriture\\des sorties};
    \node[bloc, fill=orange!20, draw=orange] (comm) at (0,-8.8)
      {\textcolor{orange}{\textbf{4.}} Communication\\\profinet{} / E/S};

    \draw[arr] (init) -- (lect);
    \draw[arr] (lect) -- (exec);
    \draw[arr] (exec) -- (ecri);
    \draw[arr] (ecri) -- (comm);

    \draw[arr, dashed, orange, very thick]
      (comm.east) -- ++(2.2,0) -- ++(0,8.8) -- (lect.east)
      node[midway, right=1.8cm, font=\small\bfseries\color{orange},
           align=center] {Cycle\\r\'ep\'et\'e\\en continu};

    \node[font=\footnotesize\itshape, color=gray, align=center] at (5.5,-4.4)
      {Dur\'ee typique\\$<$ \SI{10}{\milli\second}};
  \end{tikzpicture}
  \caption{Cycle de scrutation (scan) de l'automate \siemens{}}
  \label{fig:scan}
\end{figure}

% ============================================================
%  CHAPITRE 4 — ARCHITECTURE
% ============================================================
\chapter{Architecture G\'en\'erale du Syst\`eme}
\label{chap:archi}

\section{Architecture globale}

\begin{figure}[H]
  \centering
  \begin{tikzpicture}[
    scale=0.84, transform shape,
    box/.style={
      draw=bleuMoyen, fill=bleuClair, rounded corners=6pt,
      minimum width=3.4cm, minimum height=1.2cm,
      font=\small\bfseries, align=center,
      drop shadow={shadow xshift=2pt, shadow yshift=-2pt, opacity=0.3}
    },
    sensor/.style={box, fill=orange!20, draw=orange},
    robot/.style={box, fill=bleuFonce, text=white, draw=bleuFonce},
    arrow/.style={-{Stealth[length=7pt]}, thick},
    lbl/.style={font=\footnotesize\itshape, fill=white, inner sep=2pt}
  ]
    \node[sensor]      (capt)  at (0,4)    {Capteurs\\I0.0 -- I1.1};
    \node[sensor]      (ihm)   at (0,0)    {IHM\\Boutons + Voyants};
    \node[box]         (plc)   at (4,2)    {\api{}\\Siemens S7-1200};
    \node[box]         (krc4)  at (8,2)    {Contr\^oleur\\KRC4};
    \node[robot]       (robot) at (12,2)   {Robot\\KUKA 6 axes};
    \node[box]         (conv)  at (4,-1.5) {Convoyeur\\Moteur + VFD};
    \node[box]         (verin) at (8,-1.5) {V\'erin\\Pneumatique};

    \draw[arrow, bleuFonce]
      (capt.east) -- ++(1.2,0) |- (plc.north)
      node[pos=0.4, lbl] {\SI{24}{V} DC};
    \draw[arrow, bleuFonce]
      (ihm.east) -- ++(1.2,0) |- (plc.south)
      node[pos=0.4, lbl] {Q0.4/Q0.5};
    \draw[arrow, bleuFonce, very thick]
      (plc) -- (krc4)
      node[midway, lbl, above] {\profinet{} IO};
    \draw[arrow, bleuFonce, very thick]
      (krc4) -- (robot)
      node[midway, lbl, above] {Servo axes};
    \draw[arrow, bleuFonce]
      (plc.south) -- (conv.north)
      node[midway, lbl, right] {Q0.0};
    \draw[arrow, bleuFonce]
      (plc.east) ++(0,-0.3) -- ++(1.5,0) |- (verin.west)
      node[pos=0.3, lbl, above] {Q0.2/Q0.3};
    \draw[arrow, dashed, orange, thick]
      (krc4.south) -- (verin.north)
      node[midway, lbl, right] {\$OUT robot};

    \begin{pgfonlayer}{background}
      \node[
        draw=bleuMoyen, dashed, fill=bleuClair, fill opacity=0.15,
        rounded corners=10pt,
        fit=(plc)(krc4)(robot),
        inner sep=14pt,
        label={[bleuMoyen, font=\small\bfseries]above:\profinet{} --- 100\,Mbit/s}
      ] {};
    \end{pgfonlayer}
  \end{tikzpicture}
  \caption{Architecture globale du syst\`eme automatis\'e}
  \label{fig:archi}
\end{figure}

\section{Table des entr\'ees / sorties}

\begin{table}[H]
  \centering
  \caption{Table compl\`ete des entr\'ees/sorties de l'\api{}}
  \label{tab:es}
  \small
  \rowcolors{2}{grisClaire}{white}
  \begin{tabular}{L{2.7cm} C{1.7cm} C{1.7cm} L{6.4cm}}
    \toprule
    \rowcolor{bleuMoyen}
    \color{white}\textbf{D\'esignation} &
    \color{white}\textbf{Type} &
    \color{white}\textbf{Adresse} &
    \color{white}\textbf{Description} \\
    \midrule
    Bouton Marche    & Entr\'ee TOR & I0.0 & D\'emarrage du cycle automatique \\
    Bouton Arr\^et   & Entr\'ee TOR & I0.1 & Arr\^et en fin de cycle \\
    Capteur pi\`ece  & Entr\'ee TOR & I0.2 & D\'etection pi\`ece sur convoyeur \\
    Arr\^et d'urgence & Entr\'ee TOR & I0.3 & Arr\^et imm\'ediat s\'ecuris\'e (NF) \\
    FDC V\'erin +    & Entr\'ee TOR & I0.4 & V\'erin sorti en but\'ee \\
    FDC V\'erin --   & Entr\'ee TOR & I0.5 & V\'erin rentr\'e en but\'ee \\
    Robot\_Ready     & Entr\'ee TOR & I0.6 & Robot pr\^et \`a recevoir un ordre \\
    Robot\_End       & Entr\'ee TOR & I0.7 & Fin du cycle robot confirm\'ee \\
    Robot\_Fault     & Entr\'ee TOR & I1.0 & D\'efaut robot d\'etect\'e \\
    Mode Auto/Manuel & Entr\'ee TOR & I1.1 & S\'electeur mode de marche \\
    \midrule
    Moteur convoyeur & Sortie TOR & Q0.0 & Activation variateur convoyeur \\
    Start Robot      & Sortie TOR & Q0.1 & D\'eclenchement cycle robot \\
    V\'erin sortie   & Sortie TOR & Q0.2 & Extension v\'erin de blocage \\
    V\'erin rentr\'ee & Sortie TOR & Q0.3 & R\'etraction v\'erin \\
    Voyant Marche    & Sortie TOR & Q0.4 & Indicateur cycle actif (vert) \\
    Voyant D\'efaut  & Sortie TOR & Q0.5 & Indicateur anomalie (rouge) \\
    \bottomrule
  \end{tabular}
  \normalsize
\end{table}

% ============================================================
%  CHAPITRE 5 — ANALYSE FONCTIONNELLE
% ============================================================
\chapter{Analyse Fonctionnelle}
\label{chap:fonc}

\section{S\'equence de fonctionnement}

\begin{enumerate}[label=\textbf{\'Etape \arabic*~:}, leftmargin=*, wide, itemsep=6pt]
  \item \textbf{Initialisation} --- v\'erification de l'\'etat de tous les \'equipements
  \item \textbf{Attente pi\`ece} --- convoyeur actif, robot en position HOME
  \item \textbf{D\'etection} --- capteur I0.2 d\'etecte la pi\`ece, activation de M0.0
  \item \textbf{Arr\^et convoyeur} --- Q0.0 passe \`a 0
  \item \textbf{Blocage pi\`ece} --- Q0.2 active le v\'erin, confirmation par I0.4
  \item \textbf{Commande robot} --- Q0.1 envoie le signal Start au KRC4
  \item \textbf{Cycle robot} --- pr\'ehension et d\'epose de la pi\`ece
  \item \textbf{Fin cycle robot} --- KRC4 active \$OUT[2], r\'eception sur I0.7
  \item \textbf{Lib\'eration v\'erin} --- Q0.3 r\'etracte le v\'erin, confirmation I0.5
  \item \textbf{Red\'emarrage} --- Q0.0 repasse \`a 1, retour \'etape 2
\end{enumerate}

\section{Modes de fonctionnement}

\begin{table}[H]
  \centering
  \caption{Modes de fonctionnement du syst\`eme}
  \rowcolors{2}{grisClaire}{white}
  \begin{tabular}{L{3.5cm} L{5.5cm} L{4cm}}
    \toprule
    \rowcolor{bleuMoyen}
    \color{white}\textbf{Mode} &
    \color{white}\textbf{Description} &
    \color{white}\textbf{Activation} \\
    \midrule
    Automatique       & Cycle complet sans intervention & I1.1=1 + START \\
    Manuel            & Commande individuelle des actionneurs & I1.1=0 \\
    Pas \`a pas       & Avancement \'etape par \'etape & S\'electeur + Step \\
    D\'efaut          & Arr\^et cycle, voyant rouge & Automatique \\
    Arr\^et d'urgence & Coupure imm\'ediate s\'ecuris\'ee & E-Stop I0.3 \\
    \bottomrule
  \end{tabular}
\end{table}

% ============================================================
%  CHAPITRE 6 — GRAFCET
% ============================================================
\chapter{GRAFCET du Syst\`eme}
\label{chap:grafcet}

\section{GRAFCET de niveau 1 --- point de vue syst\`eme}

\begin{figure}[H]
  \centering
  \begin{tikzpicture}[
    scale=0.72, transform shape,
    initgrafstep/.style={
      draw=bleuFonce, fill=bleuFonce, text=white,
      minimum size=1.3cm, font=\bfseries\large,
      drop shadow={shadow xshift=2pt,shadow yshift=-2pt,opacity=0.4}
    },
    grafstep/.style={
      draw=bleuMoyen, fill=bleuClair,
      minimum size=1.3cm, font=\bfseries\large,
      drop shadow={shadow xshift=2pt,shadow yshift=-2pt,opacity=0.3}
    },
    action/.style={
      draw=orange, fill=orange!15, rounded corners=4pt,
      minimum width=6cm, minimum height=0.85cm,
      font=\small\bfseries, align=left,
      inner sep=5pt
    },
    trans/.style={font=\small\itshape, color=bleuFonce},
    arr/.style={-{Stealth[length=7pt]}, very thick, bleuFonce}
  ]

    \node[initgrafstep] (e0) at (0, 0)    {0};
    \node[grafstep]     (e1) at (0,-2.5)  {1};
    \node[grafstep]     (e2) at (0,-5)    {2};
    \node[grafstep]     (e3) at (0,-7.5)  {3};
    \node[grafstep]     (e4) at (0,-10)   {4};
    \node[grafstep]     (e5) at (0,-12.5) {5};
    \node[grafstep]     (e6) at (0,-15)   {6};
    \node[grafstep]     (e7) at (0,-17.5) {7};

    \draw[bleuFonce, very thick] (-0.85,0.85) rectangle (0.85,-0.85);

    \node[action] (a0) at (5.5, 0)
      {\textcolor{bleuFonce}{Init.} --- Reset tous bits / outputs};
    \node[action] (a1) at (5.5,-2.5)
      {\textcolor{bleuFonce}{Q0.0} = 1 \quad Convoyeur en marche};
    \node[action] (a2) at (5.5,-5)
      {\textcolor{bleuFonce}{Q0.0} = 1 \quad Attente d\'etection pi\`ece};
    \node[action] (a3) at (5.5,-7.5)
      {\textcolor{bleuFonce}{Q0.0} = 0 \quad \textcolor{orange}{Q0.2} = 1 \quad V\'erin sortie};
    \node[action] (a4) at (5.5,-10)
      {\textcolor{orange}{Q0.1} = 1 \quad Impulsion Start Robot};
    \node[action] (a5) at (5.5,-12.5)
      {\textcolor{bleuFonce}{Q0.1} = 0 \quad Attente fin cycle robot};
    \node[action] (a6) at (5.5,-15)
      {\textcolor{bleuFonce}{Q0.2} = 0 \quad \textcolor{orange}{Q0.3} = 1 \quad V\'erin rentr\'ee};
    \node[action] (a7) at (5.5,-17.5)
      {\textcolor{bleuFonce}{Q0.3} = 0 \quad \textcolor{orange}{Q0.0} = 1 \quad Red\'em.};

    \foreach \e/\a in {e0/a0,e1/a1,e2/a2,e3/a3,e4/a4,e5/a5,e6/a6,e7/a7}
      \draw[orange, thick] (\e.east) -- (\a.west);

    \foreach \s/\e/\t in {
      e0/e1/{I0.0 (START)},
      e1/e2/{Marche confirm\'ee},
      e2/e3/{I0.2 (Pi\`ece d\'etect\'ee)},
      e3/e4/{I0.4 ET Tempo 500ms},
      e4/e5/{I0.6=0 (Robot d\'emarr\'e)},
      e5/e6/{I0.7=1 (Cycle\_End)},
      e6/e7/{I0.5 (V\'erin rentr\'e)}
    }{
      \draw[arr] (\s.south) -- (\e.north);
      \path (\s.south) -- (\e.north) coordinate[pos=0.5] (mid);
      \draw[bleuMoyen, thick] ([xshift=-0.65cm]mid) -- ([xshift=0.65cm]mid);
      \node[trans, left=0.7cm of mid] {\t};
    }

    \draw[arr, dashed, orange, very thick]
      (e7.south) -- ++(0,-0.7)
      -- ++(3.5,0)
      -- ++(0,13.2)
      -- ++(-3.5,0)
      -- (e2.east);
    \node[orange, font=\small\bfseries, align=center] at (3.5,-11.5)
      {Retour\\automatique\\vers \'etape 2};

  \end{tikzpicture}
  \caption{GRAFCET de niveau 1 du syst\`eme de tri automatis\'e}
  \label{fig:grafcet}
\end{figure}

\section{GRAFCET de s\'ecurit\'e --- niveau 2}

\begin{infobox}[GRAFCET de S\'ecurit\'e --- Logique de for\c{c}age]
  \begin{itemize}
    \item \textbf{\'Etat 0} (nominal) : I0.3 = 1 (E-Stop non actionn\'e)
    \item \textbf{Transition 0$\rightarrow$1} : I0.3 = 0 (E-Stop actionn\'e, contact NF ouvert)
    \item \textbf{\'Etat 1} (urgence) : For\c{c}age du GRAFCET principal en \'etape 0 ;
          toutes sorties \`a 0 ; voyant rouge allum\'e (Q0.5)
    \item \textbf{Transition 1$\rightarrow$0} : I0.3 = 1 ET bouton d'acquittement enfonc\'e
  \end{itemize}
\end{infobox}

% ============================================================
%  CHAPITRE 7 — SCHEMAS ELECTRIQUES
% ============================================================
\chapter{Sch\'emas \'Electriques}
\label{chap:schema}

\section{Synoptique du c\^ablage g\'en\'eral}

\begin{figure}[H]
  \centering
  \begin{tikzpicture}[
    scale=0.72, transform shape,
    rail/.style={draw=bleuFonce, very thick, fill=bleuFonce!10},
    comp/.style={
      draw=bleuMoyen, fill=bleuClair, rounded corners=4pt,
      minimum width=3.6cm, minimum height=0.9cm,
      font=\small\bfseries, align=center,
      drop shadow={shadow xshift=1.5pt,shadow yshift=-1.5pt,opacity=0.3}
    },
    wire/.style={thick, bleuFonce},
    outwire/.style={thick, orange}
  ]
    \draw[rail] (-0.5,8.5) rectangle (13.5,8.9);
    \draw[rail] (-0.5,-0.4) rectangle (13.5,0);
    \node[white,font=\bfseries\small] at (0.5,8.7) {+24V};
    \node[white,font=\bfseries\small] at (0.5,-0.2) {0V};

    \node[comp, fill=orange!20, draw=orange, minimum width=4cm]
      (plc_in) at (3.5,4.5) {\api{} S7-1200\\Entr\'ees num\'eriques};

    \foreach \i/\lbl/\addr in {
      0/{Bt. START (NO)}/{I0.0},
      1/{Bt. STOP (NF)}/{I0.1},
      2/{Capteur piece PNP}/{I0.2},
      3/{E-Stop + Pilz}/{I0.3},
      4/{FDC Verin +}/{I0.4},
      5/{FDC Verin --}/{I0.5},
      6/{KUKA OUT1 Ready}/{I0.6},
      7/{KUKA OUT2 End}/{I0.7}
    }{
      \pgfmathsetmacro{\yy}{8.0 - \i * 0.85}
      \node[comp, minimum width=3.8cm, minimum height=0.7cm,
            font=\scriptsize\bfseries]
        (inp\i) at (-2.5, \yy) {\lbl};
      \draw[wire] (inp\i.east) -- (3.0, \yy)
        node[midway, above, font=\scriptsize\color{bleuFonce}] {\addr};
    }

    \node[comp, fill=bleuFonce, text=white, minimum width=4cm]
      (plc_out) at (3.5,2.5) {\api{} S7-1200\\Sorties num\'eriques};

    \foreach \i/\lbl/\addr in {
      0/{Convoyeur (KM1)}/{Q0.0},
      1/{Start Robot}/{Q0.1},
      2/{Verin sortie (YV1)}/{Q0.2},
      3/{Verin rentree (YV2)}/{Q0.3},
      4/{Voyant Marche (H1)}/{Q0.4},
      5/{Voyant Defaut (H2)}/{Q0.5}
    }{
      \pgfmathsetmacro{\yy}{3.8 - \i * 0.85}
      \node[comp, minimum width=3.8cm, minimum height=0.7cm,
            font=\scriptsize\bfseries, fill=orange!15, draw=orange]
        (out\i) at (9.5, \yy) {\lbl};
      \draw[outwire] (4.0, \yy) -- (out\i.west)
        node[midway, above, font=\scriptsize\color{orange}] {\addr};
    }
  \end{tikzpicture}
  \caption{Synoptique de c\^ablage des entr\'ees/sorties de l'\api{}}
  \label{fig:cablage}
\end{figure}

\section{Circuit de s\'ecurit\'e arr\^et d'urgence}

\begin{figure}[H]
  \centering
  \begin{tikzpicture}[
    scale=0.9, transform shape,
    relay/.style={
      draw=rougeAlerte, fill=rougeAlerte!10, rounded corners=6pt,
      minimum width=5cm, minimum height=3.5cm,
      font=\small\bfseries, align=center
    },
    contact/.style={
      draw=bleuMoyen, fill=bleuClair, rounded corners=3pt,
      minimum width=3cm, minimum height=0.7cm,
      font=\scriptsize\bfseries, align=center
    },
    arr/.style={-{Stealth[length=6pt]}, thick}
  ]
    \node[relay] (pilz) at (5,0)
      {\textbf{PILZ PNOZ X3}\\[6pt]Relais de s\'ecurit\'e\\SIL 2 / PLd};

    \node[contact, fill=rougeAlerte!20, draw=rougeAlerte]
      (estop1) at (0, 2) {E-Stop 1 (NF)};
    \node[contact, fill=rougeAlerte!20, draw=rougeAlerte]
      (estop2) at (0, 0.5) {E-Stop 2 (NF)};
    \node[font=\scriptsize, color=gray, align=center] at (0,-0.8)
      {C\^ablage bipolaire\\coupure double};

    \draw[thick, rougeAlerte]
      (estop1.east) -- ++(1,0) |- (pilz.west);
    \draw[thick, rougeAlerte]
      (estop2.east) -- ++(0.5,0) |- (pilz.west);

    \node[contact] (s1) at (10, 2.2)
      {Contact 13-14\\Circuit robot};
    \node[contact] (s2) at (10, 0.8)
      {Contact 23-24\\PLC I0.3};
    \node[contact] (s3) at (10,-0.6)
      {Contact 33-34\\Alim. sorties DO};

    \draw[arr, bleuFonce] (pilz.east) -- ++(0.5,0) |- (s1.west);
    \draw[arr, bleuFonce] (pilz.east) -- ++(0.5,0) |- (s2.west);
    \draw[arr, bleuFonce] (pilz.east) -- ++(0.5,0) |- (s3.west);

    \draw[arr, thick, bleuFonce]
      (5,3.5) -- (pilz.north)
      node[above, font=\scriptsize\bfseries] {A1: +24V / A2: 0V};

    \node[font=\scriptsize\itshape, color=gray, align=center] at (5,-2.2)
      {Homologu\'e EN ISO 13849-1 --- Cat\'egorie 3 --- PLd};
  \end{tikzpicture}
  \caption{Circuit de s\'ecurit\'e --- Relais Pilz PNOZ X3}
  \label{fig:secu_circuit}
\end{figure}

% ============================================================
%  CHAPITRE 8 — PROGRAMMATION LADDER
% ============================================================
\chapter{Programmation Ladder}
\label{chap:ladder}

\section{Symboles fondamentaux}

\begin{table}[H]
  \centering
  \caption{Symboles fondamentaux du langage Ladder (IEC 61131-3)}
  \rowcolors{2}{grisClaire}{white}
  \begin{tabular}{L{3.5cm} C{3cm} L{6cm}}
    \toprule
    \rowcolor{bleuMoyen}
    \color{white}\textbf{\'El\'ement} &
    \color{white}\textbf{Notation} &
    \color{white}\textbf{Signification} \\
    \midrule
    Contact NO        & \texttt{--[ ]--}  & Passant si bit = 1 \\
    Contact NF        & \texttt{--[/]--}  & Passant si bit = 0 \\
    Bobine            & \texttt{--( )--}  & Affectation sortie / m\'emoire \\
    Bobine Set        & \texttt{--(S)--}  & Mise \`a 1 irr\'eversible \\
    Bobine Reset      & \texttt{--(R)--}  & Mise \`a 0 irr\'eversible \\
    Front montant     & \texttt{--[P]--}  & Impulsion sur front montant \\
    Front descendant  & \texttt{--[N]--}  & Impulsion sur front descendant \\
    Temporisateur TON & \texttt{-[TON]-}  & Temporisateur \`a l'enclenchement \\
    Compteur CTU      & \texttt{-[CTU]-}  & Compteur croissant \\
    \bottomrule
  \end{tabular}
\end{table}

\section{R\'eseaux Ladder principaux}

\begin{lstlisting}[style=ladder,
  caption={Reseau 1 : Marche/Arret convoyeur avec auto-maintien}]
; Auto-maintien -- arret prioritaire -- E-Stop securise

+-[ I0.0 ]--+--[/I0.1]--[/I0.3]--[/M0.0]--(Q0.0)--+
| [START]   |  [STOP]   [ESTOP]   [Piece]           |
|           |                                        |
+-[ Q0.0 ]--+  <- auto-maintien
  [CONV]
\end{lstlisting}

\begin{lstlisting}[style=ladder,
  caption={Reseau 2 : Detection piece avec anti-rebond 50 ms}]
+-[P I0.2]--[TON T1 PT:50ms]--+
                                |
+-[ T1.Q ]---------------------( S M0.0 )-- ; Set bit piece
\end{lstlisting}

\begin{lstlisting}[style=ladder,
  caption={Reseau 3 : Commande verin de blocage}]
+-[ M0.0 ]--[/I0.3]---------------( Q0.2 )-- ; Verin sortie
+-[ M0.0 ]--[/I0.3]---------------( R Q0.3)-- ; Reset verin -
\end{lstlisting}

\begin{lstlisting}[style=ladder,
  caption={Reseau 4 : Start Robot (impulsion TP 200 ms)}]
+-[M0.0]--[I0.4]--[I0.6]--[/I0.3]--[TP T2 PT:200ms]--+
  [Piece] [Vrin+] [Ready] [EStop]                      |
+-[ T2.Q ]---------------------------------------------( Q0.1 )--
\end{lstlisting}

\begin{lstlisting}[style=ladder,
  caption={Reseau 5 : Timeout 30 s sur reponse robot}]
+-[P Q0.1]--[TON T3 PT:30000ms]--+
                                   |
+-[ T3.Q ]--[/I0.7]--------------( S M10.0 )-- ; Defaut timeout
+-[P I0.7]-------------------------( S M0.1 )-- ; Cycle_End recu
\end{lstlisting}

\begin{lstlisting}[style=ladder,
  caption={Reseau 6 : Liberation verin et redemarrage}]
+-[ M0.1 ]--[/I0.3]------------------( Q0.3 )-- ; Verin rentre
+-[ M0.1 ]--[ I0.5 ]--+--( R M0.0 )-- ; Reset piece
  [FinCyc] [Vrin-]     +--( R M0.1 )-- ; Reset Cycle_End
\end{lstlisting}

\begin{lstlisting}[style=ladder,
  caption={Reseau 7 : Gestion defauts et voyant rouge}]
+-[ I1.0  ]-------------------------( S M10.1 )-- ; Defaut robot
+-[ M10.0 ]-------------------------( S M10.1 )-- ; Defaut timeout
+-[ M10.1 ]--[/I0.3]----------------( Q0.5   )-- ; Voyant rouge
+-[ M10.1 ]--[ACQ]--[ I0.3 ]--+--( R M10.0 )--
                                +--( R M10.1 )--
\end{lstlisting}

% ============================================================
%  CHAPITRE 9 — COMMUNICATION
% ============================================================
\chapter{Communication PLC -- Robot KUKA}
\label{chap:comm}

\section{Architecture \profinet{} IO}

\begin{infobox}[Configuration r\'eseau \profinet{}]
  \begin{tabular}{@{}lll@{}}
    \textbf{\'Equipement} & \textbf{Adresse IP} & \textbf{Nom \profinet{}} \\
    \hline
    PLC S7-1200  & 192.168.1.1 & \texttt{plc-kuka} \\
    KRC4 Robot   & 192.168.1.2 & \texttt{kuka-robot} \\
  \end{tabular}\\[6pt]
  Cycle de communication : \SI{8}{\milli\second} ---
  Taille PDU : 32 octets par sens
\end{infobox}

\section{Diagramme de s\'equence PLC -- Robot}

\begin{figure}[H]
  \centering
  \begin{tikzpicture}[
    scale=0.88, transform shape,
    line/.style={very thick},
    msg/.style={-{Stealth[length=7pt]}, thick},
    lbl/.style={font=\small, fill=white, inner sep=3pt, align=center},
    time/.style={font=\scriptsize\itshape, color=gray}
  ]
    \draw[line, bleuFonce] (0,0) -- (0,-14)
      node[above, at start, font=\bfseries\normalsize, color=bleuFonce]
      {\api{} S7-1200};
    \draw[line, bleuMoyen!60!black] (8,0) -- (8,-14)
      node[above, at start, font=\bfseries\normalsize, color=bleuMoyen!60!black]
      {Robot KRC4};

    \draw[msg, bleuFonce, ->]
      (0,-1) -- (8,-1)
      node[lbl, midway, above] {Q0.1=1 : START\_ROBOT};
    \node[time] at (4,-1.5) {impulsion 200ms};

    \draw[msg, bleuMoyen!60!black, ->]
      (8,-2.5) -- (0,-2.5)
      node[lbl, midway, above] {\$OUT[1]=0 : Robot\_Ready=0};

    \draw[msg, dashed, gray, ->]
      (8,-4) -- (8,-6)
      node[lbl, midway, right=4pt] {Cycle interne robot\\(PTP, LIN, prise, depose)};

    \draw[msg, bleuMoyen!60!black, ->]
      (8,-7) -- (0,-7)
      node[lbl, midway, above] {\$OUT[2]=1 : Cycle\_End (I0.7)};

    \draw[msg, bleuFonce, ->]
      (0,-8.5) -- (8,-8.5)
      node[lbl, midway, above] {Q0.1=0 : lib\'eration START};

    \draw[msg, bleuMoyen!60!black, ->]
      (8,-10) -- (0,-10)
      node[lbl, midway, above] {\$OUT[2]=0 / \$OUT[1]=1 : Robot\_Ready};

    \draw[thick, rougeAlerte, dashed]
      (0,-1.2) -- ++(0.5,0) -- ++(0,-4)
      node[right, font=\scriptsize\bfseries\color{rougeAlerte}]
      {TON 30s} -- ++(0,-2);
    \draw[thick, rougeAlerte, ->]
      (-0.5,-7.2) -- (0,-7.2);
    \node[font=\scriptsize\color{rougeAlerte}, align=left] at (-2,-7.2)
      {Si I0.7=0\\apr\`es 30s :\\D\'efaut M10.0};

    \node[font=\small, align=left] at (4,-12.5) {%
      \textcolor{bleuFonce}{$\longrightarrow$} PLC $\rightarrow$ Robot \quad
      \textcolor{bleuMoyen!60!black}{$\longrightarrow$} Robot $\rightarrow$ PLC \quad
      \textcolor{rougeAlerte}{- - -} Surveillance timeout};

  \end{tikzpicture}
  \caption{Diagramme de s\'equence de la communication PLC -- Robot KUKA}
  \label{fig:seqdiag}
\end{figure}

\section{Table de correspondance des signaux}

\begin{table}[H]
  \centering
  \caption{Signaux PLC $\rightarrow$ Robot}
  \rowcolors{2}{grisClaire}{white}
  \small
  \begin{tabular}{C{2.2cm} C{2.1cm} L{3.4cm} L{4cm}}
    \toprule
    \rowcolor{bleuMoyen}
    \color{white}\textbf{Bit PLC} &
    \color{white}\textbf{Variable KRL} &
    \color{white}\textbf{Signal} &
    \color{white}\textbf{Description} \\
    \midrule
    Q0.1   & \texttt{\$IN[1]} & START\_ROBOT   & D\'emarrage programme robot \\
    DBX0.1 & \texttt{\$IN[2]} & RESET\_ROBOT   & R\'einitialisation apr\`es d\'efaut \\
    DBX0.2 & \texttt{\$IN[3]} & STOP\_ROBOT    & Arr\^et en fin de cycle \\
    DBX0.3 & \texttt{\$IN[4]} & PIECE\_TYPE\_A & Pi\`ece de cat\'egorie A \\
    DBX0.4 & \texttt{\$IN[5]} & PIECE\_TYPE\_B & Pi\`ece de cat\'egorie B \\
    DBX0.5 & \texttt{\$IN[6]} & VITESSE\_LENTE & Vitesse r\'eduite \\
    DBX0.6 & \texttt{\$IN[7]} & HOME\_REQ      & Demande retour HOME \\
    \bottomrule
  \end{tabular}
\end{table}
\normalsize

% ============================================================
%  CHAPITRE 10 — SECURITE INDUSTRIELLE
% ============================================================
\chapter{S\'ecurit\'e Industrielle}
\label{chap:secu}

\section{Niveaux de performance}

\begin{table}[H]
  \centering
  \caption{Fonctions de s\'ecurit\'e et niveaux de performance}
  \rowcolors{2}{grisClaire}{white}
  \begin{tabular}{L{4cm} C{2cm} C{2cm} L{5cm}}
    \toprule
    \rowcolor{bleuMoyen}
    \color{white}\textbf{Fonction} &
    \color{white}\textbf{PL requis} &
    \color{white}\textbf{SIL} &
    \color{white}\textbf{Solution} \\
    \midrule
    Arr\^et d'urgence (E-Stop) & PLd & SIL 2 & Relais Pilz PNOZ X3 \\
    Arr\^et s\'ecuris\'e robot & PLd & SIL 2 & Safe KUKA KRC4 \\
    Barri\`ere immat\'erielle  & PLe & SIL 3 & Rideau optique Sick C4000 \\
    Fin de course s\'ecuris\'e & PLc & SIL 1 & FDC \`a ouverture positive \\
    Mode manuel s\'ecuris\'e   & PLc & SIL 1 & Validation 3 positions \\
    \bottomrule
  \end{tabular}
\end{table}

\section{Zonage de s\'ecurit\'e}

\begin{figure}[H]
  \centering
  \begin{tikzpicture}[scale=0.9]
    \fill[red!20] (2,1) rectangle (7,6);
    \draw[red, very thick, dashed] (2,1) rectangle (7,6);
    \node[red, font=\bfseries\normalsize] at (4.5,6.4)
      {ZONE ROUGE --- Acc\`es interdit (mode auto)};

    \fill[orange!15] (1,0.5) rectangle (8,6.5);
    \draw[orange, thick, dashed] (1,0.5) rectangle (8,6.5);
    \node[orange, font=\bfseries\small] at (4.5,7.0)
      {ZONE ORANGE --- P\'erim\`etre surveill\'e (barri\`eres)};

    \fill[green!15] (0,0) rectangle (12,7.5);
    \draw[green!60!black, thick] (0,0) rectangle (12,7.5);
    \node[green!60!black, font=\bfseries\small] at (10.5,7.0)
      {ZONE VERTE};
    \node[green!60!black, font=\small] at (10.5,6.5) {Acc\`es libre};

    \fill[bleuFonce] (3.8,3) circle (0.7cm);
    \draw[bleuFonce, thick] (3.8,3) -- ++(0,1.5);
    \draw[bleuFonce, thick] (3.8,3) -- ++(1.2,0.8);
    \node[white, font=\bfseries\small] at (3.8,3) {R};
    \node[bleuFonce, font=\bfseries\footnotesize] at (3.8,1.5) {Robot KUKA};

    \fill[bleuMoyen!50] (5.5,2.5) rectangle (6.8,3.5);
    \node[font=\footnotesize\bfseries] at (6.15,3) {Conv.};

    \fill[green!40!black] (10,1) rectangle (11.5,2.5);
    \node[white, font=\footnotesize\bfseries, align=center] at (10.75,1.75)
      {Poste\\op\'erateur};

    \draw[orange, ultra thick, ->]
      (1.0,3) -- (2.0,3)
      node[midway, above, font=\scriptsize\bfseries\color{orange}]
      {Barri\`ere};
    \draw[orange, ultra thick, ->]
      (8.0,3) -- (7.0,3)
      node[midway, above, font=\scriptsize\bfseries\color{orange}]
      {Barri\`ere};

    \fill[red!20]        (0.2,8.2) rectangle (0.7,8.7);
    \fill[orange!15]     (2.2,8.2) rectangle (2.7,8.7);
    \fill[green!15]      (4.2,8.2) rectangle (4.7,8.7);
    \draw[red,thick]     (0.2,8.2) rectangle (0.7,8.7);
    \draw[orange,thick]  (2.2,8.2) rectangle (2.7,8.7);
    \draw[green!60!black,thick] (4.2,8.2) rectangle (4.7,8.7);
    \node[font=\scriptsize] at (1.7,8.45) {Zone interdite};
    \node[font=\scriptsize] at (3.7,8.45) {Zone surveill\'ee};
    \node[font=\scriptsize] at (5.7,8.45) {Zone libre};
  \end{tikzpicture}
  \caption{Plan de zonage de s\'ecurit\'e du poste robotis\'e}
  \label{fig:zonage}
\end{figure}

% ============================================================
%  CHAPITRE 11 — SIMULATION ET RESULTATS
% ============================================================
\chapter{Simulation et R\'esultats}
\label{chap:simu}

\section{R\'esultats des tests}

\begin{longtable}{C{0.8cm} L{4.5cm} L{3.8cm} C{2.5cm} C{1.5cm}}
  \caption{Tableau de bord des tests de simulation}
  \label{tab:tests} \\
  \toprule
  \rowcolor{bleuMoyen}
  \color{white}\textbf{N\degre} &
  \color{white}\textbf{Sc\'enario} &
  \color{white}\textbf{R\'esultat attendu} &
  \color{white}\textbf{R\'esultat} &
  \color{white}\textbf{Statut} \\
  \midrule
  \endfirsthead
  \multicolumn{5}{c}{\tablename~\thetable{} (suite)} \\
  \toprule
  \rowcolor{bleuMoyen}
  \color{white}\textbf{N\degre} &
  \color{white}\textbf{Sc\'enario} &
  \color{white}\textbf{R\'esultat attendu} &
  \color{white}\textbf{R\'esultat} &
  \color{white}\textbf{Statut} \\
  \midrule
  \endhead
  \rowcolor{grisClaire}
  T01 & D\'emarrage syst\`eme normal     & Robot HOME, convoyeur actif         & Conforme   & \textcolor{green!60!black}{\checkmark} \\
  T02 & D\'etection pi\`ece + arr\^et    & Convoyeur stop $<$ 50ms             & 45ms       & \textcolor{green!60!black}{\checkmark} \\
  \rowcolor{grisClaire}
  T03 & Cycle robot pi\`ece type A       & D\'epose zone A en 8.2s             & 8.1s       & \textcolor{green!60!black}{\checkmark} \\
  T04 & Cycle robot pi\`ece type B       & D\'epose zone B en 9.5s             & 9.4s       & \textcolor{green!60!black}{\checkmark} \\
  \rowcolor{grisClaire}
  T05 & E-Stop en cours de cycle         & Arr\^et $<$ 200ms s\'ecuris\'e      & 180ms      & \textcolor{green!60!black}{\checkmark} \\
  T06 & D\'efaut robot simul\'e          & Voyant rouge, arr\^et cycle          & Conforme   & \textcolor{green!60!black}{\checkmark} \\
  \rowcolor{grisClaire}
  T07 & Timeout communication robot      & D\'efaut apr\`es 30s                & 29.8s      & \textcolor{green!60!black}{\checkmark} \\
  T08 & Red\'emarrage apr\`es E-Stop     & Retour cycle apr\`es acquittement   & Conforme   & \textcolor{green!60!black}{\checkmark} \\
  \rowcolor{grisClaire}
  T09 & Mode pas \`a pas                 & Avancement par \'etapes             & Conforme   & \textcolor{green!60!black}{\checkmark} \\
  T10 & 100 cycles cons\'ecutifs         & 0 erreur, cadence stable            & 42s/cyc    & \textcolor{green!60!black}{\checkmark} \\
  \bottomrule
\end{longtable}

\section{Indicateurs de performance}

\begin{table}[H]
  \centering
  \caption{Performances mesur\'ees lors de la simulation}
  \rowcolors{2}{grisClaire}{white}
  \begin{tabular}{L{5cm} C{3cm} C{3cm} C{2cm}}
    \toprule
    \rowcolor{bleuMoyen}
    \color{white}\textbf{Indicateur} &
    \color{white}\textbf{Mesur\'e} &
    \color{white}\textbf{Objectif} &
    \color{white}\textbf{\'Ecart} \\
    \midrule
    Cadence de production       & 85 pi\`eces/h & 80 pi\`eces/h     & \textcolor{green!60!black}{+6.3\%} \\
    Temps de cycle moyen        & \SI{42.3}{s}  & $<$ \SI{45}{s}    & \textcolor{green!60!black}{OK} \\
    Temps de r\'eponse E-Stop   & \SI{180}{ms}  & $<$ \SI{200}{ms}  & \textcolor{green!60!black}{OK} \\
    Taux de d\'etection pi\`ece & 100\%         & $\geq$ 99.5\%     & \textcolor{green!60!black}{OK} \\
    Taux d'erreur position.     & 0.02\%        & $<$ 0.1\%         & \textcolor{green!60!black}{OK} \\
    Disponibilit\'e (8 heures)  & 99.8\%        & $>$ 99\%          & \textcolor{green!60!black}{OK} \\
    \bottomrule
  \end{tabular}
\end{table}

% ============================================================
%  CHAPITRE 12 — CONCLUSION
% ============================================================
\chapter{Conclusion G\'en\'erale}

\section{Bilan du projet}

Ce projet d'automatisme industriel a permis de mettre en pratique l'ensemble des
comp\'etences acquises en g\'enie \'electrique et en automatique industrielle. La conception
et la r\'ealisation du syst\`eme de tri automatis\'e int\'egrant un robot \kuka{} et un
automate \siemens{} constituent un cas industriel complet.

\begin{itemize}[itemsep=4pt]
  \item[$\checkmark$] Architecture mat\'erielle et logicielle d\'efinie et document\'ee
  \item[$\checkmark$] Programmation Ladder compl\`ete sous \tia{}
  \item[$\checkmark$] Communication \profinet{} PLC--Robot op\'erationnelle
  \item[$\checkmark$] GRAFCET mod\'elisant le comportement s\'equentiel complet
  \item[$\checkmark$] Fonctions de s\'ecurit\'e PLd/SIL2 impl\'ement\'ees
  \item[$\checkmark$] Simulation valid\'ee, performances d\'epassant les objectifs
\end{itemize}

\section{Perspectives d'am\'elioration}

\begin{enumerate}[label=\textbf{\arabic*.}, itemsep=4pt]
  \item Int\'egration d'un syst\`eme de vision industrielle pour le tri par forme et couleur
  \item D\'eploiement d'un SCADA pour la supervision et la tra\c{c}abilit\'e
  \item Maintenance pr\'edictive bas\'ee sur l'analyse des donn\'ees de cycle (IIoT)
  \item Extension \`a une architecture multi-robots en r\'eseau
  \item Certification ISO 10218 pour robot collaboratif (cobot)
\end{enumerate}

% ── Bibliographie ────────────────────────────────────────────
\nocite{*}
\printbibliography[heading=bibintoc, title={Bibliographie}]

% ── Annexes ──────────────────────────────────────────────────
\appendix
\chapter{Param\`etres de configuration TIA Portal}

\begin{infobox}[Configuration mat\'erielle TIA Portal V17]
  \begin{itemize}
    \item CPU : S7-1200 CPU 1214C DC/DC/DC --- R\'ef. : 6ES7 214-1AG40-0XB0
    \item Adresse IP : 192.168.1.1 --- Masque : 255.255.255.0
    \item Nom \profinet{} : \texttt{plc-kuka}
    \item Cycle \profinet{} : \SI{8}{\milli\second}
    \item Mode : RUN-P
  \end{itemize}
\end{infobox}

\chapter{R\'ef\'erences normatives}

\begin{itemize}
  \item IEC 61131-3:2013 --- Langages de programmation des automates
  \item IEC 60848:2013 --- Langage GRAFCET
  \item EN ISO 13849-1:2015 --- S\'ecurit\'e des machines, parties SCS
  \item IEC 62061:2021 --- S\'ecurit\'e fonctionnelle
  \item IEC 61158 --- R\'eseaux de terrain industriels (\profinet{})
  \item ISO 10218-2:2011 --- Robots industriels
  \item Directive 2006/42/CE --- Directive Machines
\end{itemize}

\end{document}
